\documentclass{article}

\input{setup}

\begin{document}

\CAPA{TP1 - Problema do Caixeiro Viajante}{BCC202 - Estrutura de Dados I}{Gustavo Ferreira e Marcelo Vieira}{Pedro Silva}

\section{Introdução}
O problema do caixeiro viajante é antigo na Ciência da Computação, e consiste em determinar a melhor rota para percorrer uma série de cidades, retornando à cidade de origem. Para realizar esse trabalho, foram desenvolvido um algoritmo em C e um relatório referente ao que foi implementado. A codificação foi feita utilizando somente a biblioteca padrão da GNU, sem a utilização de bibliotecas adicionais.

\subsection{Especificações do problema}
Para resolver o problema do caixeiro viajante, foi definido previamente que uma função recursiva deve ser utilizada a fim de encontrar a menor distância possível entre as cidades. Além disso, um Tipo Abstrato de Dados (TAD) chamado GrafoPonderado também precisou ser implementado como forma de representar o caminho a ser analisado, ele possui dois atributos pré definidos: número de cidades a serem visitadas e uma matriz de adjacências. O TAD também deve implementar pelo menos 5 operações: 
\begin{enumerate}
    \item alocarGrafo: aloca um (ou mais) TAD GrafoPonderado.
    \item desalocarGrafo: desaloca um TAD GrafoPonderado.
    \item leGrafo: inicializa o TAD GrafoPonderado a partir de dados do terminal.
    \item encontraCaminho: função recursiva que retorna o menor caminho no grafo fornecido.
    \item imprimeCaminho: imprime na tela o menor caminho e a distância total percorrida.
\end{enumerate}


\subsection{Considerações iniciais}

Algumas ferramentas foram utilizadas durante a criação deste projeto:

\begin{itemize}
  \item Ambiente de desenvolvimento do código fonte: Visual Studio Code. \footnote{Visual Studio Code está disponível em \url{https://code.visualstudio.com/}}
  \item Linguagem utilizada: C.
  \item Ambiente de desenvolvimento da documentação: Overleaf \LaTeX. \footnote{Disponível em \url{https://www.overleaf.com/}}
\end{itemize}

\subsection{Ferramentas utilizadas}
Algumas ferramentas foram utilizadas para testar a implementação, como:
\begin{itemize}
    \item[-] \textit{Valgrind}: ferramentas de análise dinâmica do código.
\end{itemize}

\subsection{Especificações da máquina}
A máquina onde o desenvolvimento e os testes foram realizados possui a seguinte configuração:
\begin{itemize}
    \item[-] Processador: Ryzen 5 5600G.
    \item[-] Memória RAM: 16Gb.
    \item[-] Sistema Operacional: Linux.
\end{itemize}

\subsection{Instruções de compilação e execução}

Para a compilação do projeto, basta digitar:

\begin{tcolorbox}[title=Compilando o projeto,width=\linewidth]
    gcc -g -o exe arquivo1.c arquivo2.c -Wall
\end{tcolorbox}

Usou-se para a compilação as seguintes opções:
\begin{itemize}
    \item [-] \emph{-o exe}: para definir o nome do arquivo executável gerado.
    \item [-] \emph{-g}: para compilar com informação de depuração e ser usado pelo Valgrind.
    \item [-] \emph{-Wall}: para mostrar todos os possível \emph{warnings} do código.
\end{itemize}

Para a execução do programa basta digitar:
\begin{tcolorbox}[title=,width=\linewidth]
    ./exe caminho\_até\_o\_arquivo\_de\_entrada
\end{tcolorbox}


\clearpage
\section{Desenvolvimento}
Nessa seção serão descritos os passos realizados para a implementação do programa, incluindo as dificuldades encontradas e as soluções definidas. 

\subsection{O Tipo Abstrato Definido}
A fim de atender aos requisitos impostos, um tipo de dados chamado GrafoPon foi definido. O Código \ref{lst:cod1} apresenta o código implementado para a struct GrafoPon em C, contendo os atributos e suas devidas explicações em forma de comentários.

\begin{lstlisting}[caption={Implementação da struct GrafoPon.},label={lst:cod1},language=C]
typedef struct grafoPon {
    int numCidades; //numero de cidades totais
    int **matrizAdjacencias; //matriz que armazena as distancias entre cada cidade
    int *cidadeVisitada; //vetor de 0 e 1 que nos diz se uma cidade foi visitada ou nao
    int *caminho; //vetor que armazena o melhor caminho pro dado grafo, uma vez que calculado
    int *caminhoTemp; //vetor que armazena os caminhos em tempo de teste
    int distancia; //inteiro que armazena a melhor distancia encontrada
} GrafoPon;
 \end{lstlisting}
\subsection{Alocando memória e lendo os dados}
Com o tipo abstrato definido, o próximo passo consiste em alocar dinâmicamente a memória para os atributos necessários e realizar a leitura dos dados. Para isso, foram criadas duas funções: "leGrafo" e "alocarGrafo". A função para ler os grafos primeiramente recebe o número de cidades que o caixeiro viajante precisará visitar, e logo em seguida chama a função responsável por alocar espaço para que a matriz de ajdacências seja devidamente alocada. Essa matriz será de tamanho n X n, onde n representa o número de cidades a serem visitadas. A função também aloca espaço para os três vetores da struct GrafoPon: cidadeVisitada, caminhoTemp e caminho. A melhor distância encontrada também é definida utilizando o parâmetro INT\_MAX, da biblioteca limits.h, que define a melhor distância como sendo o maior número inteiro permitido pelo sistema.
Após o espaço ser alocado o algoritmo volta para a função de leitura, que lê e armazena as distâncias entre as cidades no seguinte padrão: sendo i uma cidade A e j uma cidade B, um elemento [i][j] da matriz representa a distância entre as cidades A e B.

Logo após a implementação das funções para leitura de dados e alocação de memória, também foi implementada a função "desalocarGrafo", que é chamada ao final da execução do algoritmo e responsável por desalocar a memória de todos os atributos que foram alocados dinâmicamente na função "alocarGrafo".

Com isso, três das operações pré definidas já se encontravam implementadas.

\subsection{Imprimindo os resultados}
Após a implementação de três das cinco funções necessárias para o funcionamento do programa, foi decidido que o próximo passo seria implementar a função que imprime os resultados. Por ser uma função simples cujos parâmetros já são conhecidos, foi possível realizar sua implementação antes mesmo de definir a função que encontra o melhor caminho. O código \ref{lst:cod1} apresenta a implementação da função "imprimeCaminho".
\begin{lstlisting}[caption={Implementação da função imprimeCaminho.},label={lst:cod1},language=C]
void imprimeCaminho(GrafoPon *grafo) {
    //imprime numa primeira linha o menor caminho a ser percorrido
    for (int i = 0; i < grafo->numCidades; i++) {
        printf("%d ", grafo->caminho[i]);
    }
    //como a ultima cidade sempre sera o 0, imprimimos ele ao final
    printf("0 ");

    //imprime numa segunda linha a distancia percorrida por esse menor caminho
    printf("\n%d\n", grafo->distancia);
}
 \end{lstlisting}

\subsection{Encontrando o melhor caminho}
Com quatro das cinco operações implementadas, partimos para a função "encontraCaminho", responsável por encontrar o menor caminho possível entre as cidades na qual o caixeiro viajante deve passar. Inicialmente, foi utilizado um algoritmo de heap dentro de um laço for que realizava todas as permutações de n objetos, porém, o algoritmo não se mostrou a melhor saída quando combinado à estrutura definida e nem todas as permutações estavam sendo realizadas. Assim, a fim de conseguir gerar todas as permutações dentro do laço for de maneira simplificada, foram utilizados parâmetros inteiros para substituir as permutações por heap.

O código \ref{lst:cod1} apresenta o protótipo da função encontraCaminho.
\begin{lstlisting}[caption={Protótipo da função encontraCaminho.},label={lst:cod1},language=C]
void encontraCaminho(GrafoPon *grafo, int posicao, int contador, int distanciaTemp);
 \end{lstlisting}

 O primeiro parâmetro que a função recebe é um ponteiro que apontada para o grafo do algoritmo, ele será importante para termos acesso ao número de cidades e à matriz de ajdacências. Além disso, com acesso direto ao grafo conseguimos atualizar a qualquer momento a melhor distância, o vetor contendo o melhor caminho, o vetor contendo os caminhos temporários definidos e o vetor que informa qual cidade foi visitada. O parâmetro posição informa a cidade atual que o caixeiro viajante está visitando, e como ele parte da primeira cidade, sempre iremos iniciar a função passando o número 0 como parâmetro. O contador é o parâmetro que conta o número de cidades visitadas até o momento, ele deve ser inicializado como 1, pois o caixeiro já visitou a primeira cidade. O último parâmetro, chamado de distanciaTemp, irá armazenar a distância calculada atualmente, que começa em 0 e vai sendo incrementado conforme o caixeiro visita outras cidades.
 
 O código \ref{lst:cod1} apresenta a implementação completa da função encontraCaminho e como os parâmetros são utilizados.
 \begin{lstlisting}[caption={Implementação da função encontraCaminho.},label={lst:cod1},language=C]
void encontraCaminho(GrafoPon *grafo, int posicao, int contador, int distanciaTemp) {
    if ((contador == grafo->numCidades) && (grafo->matrizAdjacencias[posicao][0])) {
        distanciaTemp += grafo->matrizAdjacencias[posicao][0];
        if (distanciaTemp < grafo->distancia) {
            grafo->distancia = distanciaTemp;
            for (int k = 0; k < grafo->numCidades; k ++)
                grafo->caminho[k] = grafo->caminhoTemp[k];
        }
        return;
    }    
    for (int i = 0; i < grafo->numCidades; i++) {
        if (grafo->cidadeVisitada[i] == 0 && grafo->matrizAdjacencias[posicao][i]) {
            grafo->cidadeVisitada[i] = 1;
            grafo->caminhoTemp[contador] = i;
            encontraCaminho(grafo, i, contador + 1, distanciaTemp + grafo->matrizAdjacencias[posicao][i]);
            grafo->cidadeVisitada[i] = 0;
        }
    }
}
 \end{lstlisting}

 Primeiramente, a função irá conferir se todas as cidades já foram visitadas pelo caixeiro viajante no caminho atual e se a distância entre a cidade atual e a inicial é diferente de 0. O algoritmo inicialmente desenvolvido apenas conferia se todas as cidades haviam sido visitadas, não conferindo se a distância entre a cidade atual e inicial era diferente de 0, porém, em determinados casos de teste essa condição se tornou necessária, pois algumas vezes determinadas cidades não possuem ligação com a cidade inicial.

 Caso as duas condições sejam satisfeitas, a distância total do caminho atual é calculada, e caso ela seja menor do que a menor distância encontrada até o momento, atualizamos a menor distância para receber a distância do caminho atual. Além de atualizar a menor distância, também é necessário atualizar o vetor que armazena a ordem das cidades do caminho para que receba a ordem das cidades do caminho atual.

 O laço for irá iterar sobre todas as cidades, caso a cidade da iteração não tenha sido visitada e haja uma distância diferente de 0 (o que significa ser possível viajar entre elas) entre ela e a cidade atual, a cidade da iteração é visitada e seu número é inserido no vetor que armazena o caminho temporário. A função é então chamada novamente para que a próxima cidade seja visitada, com parâmetros atualizados:
 \begin{enumerate}
    \item Posição recebe o número da iteração, que representa a próxima cidade a ser visitada;
    \item O contador é incrementado, pois o caixeiro acaba de visitar uma cidade;
    \item A distância do caminho atual é incrementada, sendo somada com a distância entre a cidade atual e a cidade da iteração;
\end{enumerate}
Após a chamada recursiva, a cidade deve ser marcada como não visitada para que as próximas permutações possam ser realizadas normalmente.

\subsection{Implementação da função principal}
Com todo o TAD definido, resta criar um grafo na função principal e realizar as operações na ordem correta. O código \ref{lst:cod1} apresenta a implementação da função main e as chamadas de funções.
 \begin{lstlisting}[caption={Implementação da função main.},label={lst:cod1},language=C]
int main() {

    GrafoPon grafo;

    leGrafo(&grafo);

    encontraCaminho(&grafo, 0, 1, 0);

    imprimeCaminho(&grafo);

    desalocarGrafo(&grafo);
    
    return 0;
}
 \end{lstlisting}

\clearpage

\section{Experimetos}
Após implementação do código, uma série de experimentos foram feitos testando diferentes entradas. Os números de cidades foram testados conforme arquivos de testes enviados, utilizando a IDE Visual Studio Code e o terminal do Linux. O arquivo de testes é composto de 5 entradas, e o número de cidades definido como 4, 6, 4, 5 e 5. Dois dos cinco arquivos possuíam cidades cuja distância uma para a outra era de 0, ou seja, não há como o caixeiro viajante ir de uma para a outra. O tempo de execução do arquivo de testes também foi testado.
\clearpage

\section{Resultados}

O algoritmo se mostrou eficiente para a resolução de todos os casos de testes, conforme mostrado na Figura 1.
\begin{figure}[!htb]
    \centering
    \includegraphics[width=0.7\textwidth]{testesfeitos.jpeg}
    \caption{Resultado da comparação entre saídas + o tempo de execução do programa.}
    \label{fig:testes}
\end{figure}

\clearpage

\section{Considerações Finais}
Após a realização do trabalho, foi possível compreender melhor sobre a utilização de alocação dinâmica, structs e funções recursivas. Também foi possível aprender sobre processos para debugar um programa, tendo em vista que durante o desenvolvimento do algoritmo vários bugs foram encontrados.

\clearpage

\end{document}